\documentclass[11pt]{article}
\usepackage{amsmath,amssymb,amsthm}
\usepackage[T1]{fontenc}
\usepackage{lmodern}
\usepackage[margin=1in]{geometry}
\usepackage{amsmath,amssymb,amsthm,mathtools,mathrsfs}
\usepackage[hidelinks]{hyperref}
\newtheorem{theorem}{Theorem}[section]
\newtheorem{proposition}[theorem]{Proposition}
\newtheorem{lemma}[theorem]{Lemma}
\newtheorem{corollary}[theorem]{Corollary}
\theoremstyle{definition}
\newtheorem{definition}[theorem]{Definition}
\theoremstyle{remark}
\newtheorem{remark}[theorem]{Remark}
\newcommand{\kk}{\mathbb{K}}
\newcommand{\Z}{\mathbb{Z}}
\newcommand{\Q}{\mathbb{Q}}
\newcommand{\eps}{\varepsilon}
\newcommand{\id}{\operatorname{id}}
\newcommand{\sgnind}[1]{\mathbf{1}_{\{#1\}}}
\newcommand{\cA}{\mathcal{A}}
\newcommand{\cT}{\mathcal{T}}
\newcommand{\cD}{\mathcal{D}}
\newcommand{\cH}{\mathcal{H}}
\newcommand{\cS}{\overline{\mathscr{S}}}

\newcommand{\fptitle}{Review of Problem 3 -- MerLean}
\newcommand{\reviewtag}{Reviewer 1}
\usepackage{fp-review}

\begin{document}
\maketitle

\section{Review}
The solution only claims a partial result. The solution engages with the algebraic part of the problem, and acknowledges that it does not attempt to engage with the skein theoretic portion. 

\paragraph{Correctness.}
The paper acknowledges a partial result, and that partial result is the easier half of the main problem statement. But what the partial solution accomplished does appear to be correct.

\paragraph{Exposition.}
I had trouble finding the precise statement that serves as the objective of the partial result. In each section, I found it hard to deduce what goal was accomplished.

\paragraph{Attribution.}

There are too few papers cited. The main source for stated $SL_n$ skein algebras should be L\^{e}-Sikora, which is not included in the bibliography. 

\paragraph{Correctness.}
\begin{itemize}
	\item  1: essentially flawless mathematical argument
	\item  0.5: minor errors that are easily corrected
	\item  0: blank solution, or acknowledged partial result
	\item  $-0.5$: major errors
	\item  $-1$: proposed proof makes no progress towards the result.
\end{itemize}
Note that points should be docked for wrong statements/ solutions.

\noindent Score: 0


\paragraph{Exposition.}
\begin{itemize}
	\item  1: well-written and clear solution.
	\item  0: argument that is possible to follow with serious effort
	\item  $-1$: very poor exposition
\end{itemize}

\noindent Score: 0


\paragraph{Attributions.}
\begin{itemize}
	\item  1: solutions contains fully appropriate attributions
		to related work and theorems being used.
	\item  0: moderate issues with missing attributions or hallucinated references
	\item  $-1$: blatant plagiarism or hallucinations
\end{itemize}

\noindent Score: 0


\newpage
\section{Submission with inline comments}


\begin{abstract}
The problem considered here asks for an explicit identification of a reduced stated skein algebra of a punctured bigon with a generalized double of the full-weight quantum Borel algebras, followed by a diagonal quotient giving the simply-connected quantum group. We give a partial answer. At a specified generic parameter we compute the ordered double, its full diagonal quotient, and explicit inverse formulas for triangular quantum matrices. These yield a comparison theorem for an algebra satisfying a precise mixed exchange identity and a vector-space factorization. The proof retains the weight lattice and distinguishes gluing all weight Cartans from gluing only simple-root Cartans. The identification of these hypotheses with the designated skein algebra and its arc conventions remains unresolved; consequently neither requested isomorphism for the original objects is asserted.
\end{abstract}

\section{The question and the scope of the reduction}

Let $U_q(\mathfrak b_+)$ and $U_q(\mathfrak b_-)$ denote the upper and lower Borel subalgebras of the simply-connected quantum group $U_q(\mathfrak{sl}_n)$. The question is to construct
\begin{equation}\label{eq:question}
 \cS_n(\mathfrak B^*)\ \cong\ U_q(\mathfrak b_-)\bowtie U_q(\mathfrak b_+)
\end{equation}
for the prescribed Drinfeld--Tanisaki pairing
\begin{equation}\label{eq:requested-pairing}
 \begin{gathered}
 \langle\hat K_r,\hat K_u\rangle=q^{1/n-\delta_{ru}},\qquad
 \langle E_i,F_j\rangle=\frac{\delta_{ij}}{q^{-1}-q},\\
 \langle\hat K_r,F_j\rangle=\langle E_i,\hat K_u\rangle=0,
 \end{gathered}
\end{equation}
and then to deduce
\begin{equation}\label{eq:question-quotient}
 U_q(\mathfrak{sl}_n)\ \cong\ \cS_n(\mathfrak B^*)/
 (b_{rr}c_{rr}-1:1\leq r\leq n).
\end{equation}
Here $b_{ij}$ and $c_{ij}$ are the two designated families of rightward horizontal stated arcs. A state is an endpoint label in $\{1,\ldots,n\}$. Parentheses in an algebra quotient always denote the two-sided ideal generated by the listed elements.

An explicit solution requires more than a matching commutator. The endpoint order, boundary reduction and multiplication convention on the skein side must determine the same mixed relation as the ordered double. Injectivity must then be proved. Furthermore, the simply-connected form involves weight Cartans, not just simple-root Cartans.

The precise definition of the original $\cS_n(\mathfrak B^*)$, its designated arc diagram, and its coefficient and fractional-power conventions are not fixed here. We therefore keep \eqref{eq:question}--\eqref{eq:question-quotient} as the original targets and distinguish them from all algebras defined below. The results of this paper are a generic algebraic reduction of these targets. The missing geometric and convention identifications are stated in Section~\ref{sec:application}.

Our main comparison result is Theorem~\ref{thm:comparison}. It gives every triangular matrix entry as a root vector with its exact scalar and Cartan factor. Its hypotheses include an explicit mixed exchange identity and a multiplication factorization; it does not presume the desired isomorphism. The independent presentation calculation in Theorem~\ref{thm:double} identifies the resulting algebra with a specified ordered double. Proposition~\ref{prop:quotient} proves the diagonal quotient, and Proposition~\ref{prop:strict} explains why gluing only simple-root ratios is insufficient.

The full-weight presentations and pairing used here are recalled from Voigt--Yuncken \cite[Definition 2.13, Proposition 2.14, and Section 2.8.1]{VY}; we explicitly convert their coproduct convention. The inverse triangular-matrix calculation uses the type-A root relations recorded by Chen--Shao--Shum \cite[Section 3]{CSS}, who attribute those relations to Yamane. We use the relations, rather than their basis theorem. On the geometric side, L\^e--Yu \cite[Theorems 6.2 and 7.2]{LY} provide arc factorization and normality of the ideals generated by bad corner arcs. Their Theorem 6.2 is credited to L\^e--Sikora. \review{L\^{e}--Sikora is mentioned, but should be cited in the bibliography.} Section~\ref{sec:application} explains how such results could supply a factorization after reduction, without identifying the original diagram by assumption. No novelty claim is made for the underlying quantum-group theory.

\section{Weight Cartans and the two presentations}\label{sec:presentations}

Throughout the algebraic results, fix an integer $n\geq2$, an indeterminate $t$, and
\begin{equation}\label{eq:base}
 \kk=\Q(t),\qquad q=t^{2n^2},\qquad d=q-q^{-1},\qquad q^{1/n}=t^{2n}.
\end{equation}
Thus $q$ is not a root of unity and $d\ne0$. All tensor products are over $\kk$ unless another base is specified. We use $1\leq r,u\leq n$ for weight indices and $1\leq i,j<n$ for simple-root indices. Set
\[
 w_{ri}=\delta_{ri}-\delta_{r,i+1},\qquad
 a_{ij}=2\delta_{ij}-\delta_{i,j+1}-\delta_{i+1,j}.
\]

\subsection{Why the determinant-one torus is the weight torus}
Let $\eps_1,\ldots,\eps_n$ be the standard basis of $\Q^n$ and put
\[
 \lambda_r=\eps_r-\frac1n\sum_{u=1}^n\eps_u,\qquad
 W=\{x\in\Q^n:\textstyle\sum_r x_r=0,\ x_r-x_u\in\Z\text{ for all }r,u\}.
\]
The lattice $W$, with the ordinary dot product, is the type-$A_{n-1}$ weight lattice. Its simple roots are $\alpha_i=\lambda_i-\lambda_{i+1}$.

\begin{lemma}\label{lem:lattice}
The map $[m]\mapsto\sum_r m_r\lambda_r$ identifies $\Z^n/\Z(1,\ldots,1)$ with $W$. The fundamental weights are $\omega_i=\lambda_1+\cdots+\lambda_i$. Moreover,
\[
 (\lambda_r,\lambda_u)=\delta_{ru}-1/n,\qquad
 (\lambda_r,\alpha_i)=w_{ri}.
\]
The least positive integer $L$ for which $(W,W)\subseteq L^{-1}\Z$ is $n$.
\end{lemma}
\begin{proof}
The image of $m$ has coordinates $m_r-\frac1n\sum_u m_u$. It is zero precisely when $m$ is constant. Conversely, for $x\in W$ take the integers $m_r=x_r-x_n$; their image is $x$. The displayed inner products follow by expansion. They give $(\omega_i,\alpha_j)=\delta_{ij}$, while
$\lambda_1=\omega_1$, $\lambda_r=\omega_r-\omega_{r-1}$ for $1<r<n$, and $\lambda_n=-\omega_{n-1}$. Thus these are all the fundamental weights and they generate $W$. For integer representatives $m,l$,
\[
 \left(\sum_r m_r\lambda_r,\sum_r l_r\lambda_r\right)
 =\sum_r m_rl_r-\frac{(\sum_r m_r)(\sum_r l_r)}n.
\]
This belongs to $n^{-1}\Z$, and $(\lambda_1,\lambda_1)=(n-1)/n$ has denominator $n$.
\end{proof}

\subsection{Borel and double presentations}
Define $B_+$ to be the algebra generated by commuting units $P_r$ and elements $E_i$, subject to
\begin{equation}\label{eq:positive}
 \prod_rP_r=1,\qquad P_rE_iP_r^{-1}=q^{w_{ri}}E_i.
\end{equation}
The root generators satisfy the type-A quantum Serre relations: for $X_i=E_i$,
\begin{equation}\label{eq:serre}
 \begin{aligned}
 X_iX_j&=X_jX_i &&(|i-j|>1),\\
 X_i^2X_j-(q+q^{-1})X_iX_jX_i+X_jX_i^2&=0 &&(|i-j|=1).
 \end{aligned}
\end{equation}
Define $B_-$ by commuting units $N_r$ with $\prod_rN_r=1$, generators $F_i$, the same root relations, and
\begin{equation}\label{eq:negative}
 N_rF_iN_r^{-1}=q^{-w_{ri}}F_i.
\end{equation}
For $\lambda=\sum m_r\lambda_r$ write $P_\lambda=\prod P_r^{m_r}$ and $N_\lambda=\prod N_r^{m_r}$. Lemma~\ref{lem:lattice} and the product-one relations make these well defined. Abbreviate
\[
 A_i=P_iP_{i+1}^{-1},\qquad B_i=N_iN_{i+1}^{-1}.
\]
Here the indexed element $B_i$ is distinct from the algebra $B_-$.

Define $\cA$ by the two Borel presentations together with
\begin{equation}\label{eq:mixed}
 \begin{aligned}
 P_rN_u&=N_uP_r,& P_rF_i&=q^{-w_{ri}}F_iP_r,\\
 N_rE_i&=q^{w_{ri}}E_iN_r,&
 [E_i,F_j]&=\delta_{ij}\frac{A_i-B_i^{-1}}d.
 \end{aligned}
\end{equation}
As usual, $[x,y]=xy-yx$. Define $\cT$ to have one family of commuting units $K_r$ of product one, the roots $E_i,F_i$, their relations \eqref{eq:serre}, weights $q^{w_{ri}}$ and $q^{-w_{ri}}$, and
\begin{equation}\label{eq:single-commutator}
 [E_i,F_j]=\delta_{ij}\frac{L_i-L_i^{-1}}d,
 \qquad L_i=K_iK_{i+1}^{-1}.
\end{equation}
These are definitions of auxiliary algebras, independent of any skein identification.

\section{The pairing and the ordered double}\label{sec:double}

We first realize $B_\pm$ as actual Hopf algebras. A Hopf algebra has a coproduct $\Delta$, a counit $\epsilon$, and an antipode $S$. We write $\Delta(x)=\sum x_1\otimes x_2$ and $(\Delta\otimes\id)\Delta(x)=\sum x_1\otimes x_2\otimes x_3$, suppressing summation indices. All such sums below are finite for a fixed element. A group-like element $g$ satisfies $\Delta(g)=g\otimes g$ and $\epsilon(g)=1$; its antipode is $g^{-1}$.

\begin{proposition}\label{prop:pairing}
The presentations $B_\pm$ admit Hopf structures with group-like Cartans and
\begin{equation}\label{eq:hopf}
 \begin{aligned}
 \Delta(E_i)&=E_i\otimes1+A_i\otimes E_i,&S(E_i)&=-A_i^{-1}E_i,\\
 \Delta(F_i)&=F_i\otimes B_i^{-1}+1\otimes F_i,&S(F_i)&=-F_iB_i.
 \end{aligned}
\end{equation}
The Cartan counits are one and the root counits are zero. There is a bilinear pairing $\tau:B_+\times B_-\to\kk$ satisfying
\begin{equation}\label{eq:skew}
 \begin{aligned}
 \tau(xx',y)&=\sum\tau(x',y_1)\tau(x,y_2),\\
 \tau(x,yy')&=\sum\tau(x_1,y)\tau(x_2,y'),\\
 \tau(1,y)&=\epsilon(y),\qquad \tau(x,1)=\epsilon(x),
 \end{aligned}
\end{equation}
and
\begin{equation}\label{eq:table}
 \begin{gathered}
 \tau(P_r,N_u)=q^{1/n-\delta_{ru}},\qquad
 \tau(E_i,F_j)=-\delta_{ij}/d,\\
 \tau(P_r,F_j)=\tau(E_i,N_u)=0.
 \end{gathered}
\end{equation}
\end{proposition}
\begin{proof}
Use the full-weight algebra of \cite[Definition 2.13]{VY} for type $A_{n-1}$. The source rank is $n-1$, its squared root lengths are two, and its integer $L$ is $n$ by Lemma~\ref{lem:lattice}. Its parameter $s$ can be taken to be $t^{2n}$, so its $q=s^L$ agrees with \eqref{eq:base}. This checks both the fractional-power convention and the non-root-of-unity hypothesis used for its pairing.

The paragraph after \cite[Proposition 2.14]{VY} identifies the actual Borel subalgebras with their corresponding presentations. Send their $K_\lambda$ to $P_\lambda$ or $N_\lambda$, respectively, and fix the root generators. The inverse assignments are $P_r\mapsto K_{\lambda_r}$ and $N_r\mapsto K_{\lambda_r}$. Lemma~\ref{lem:lattice} checks the torus relations; the weights are $(\lambda_r,\alpha_i)=w_{ri}$, and the type-A Serre relations are exactly \eqref{eq:serre}. These are inverse presentation maps on the actual Borels.

Write $\Delta_0,S_0$ for the transported source Hopf operations. Their root formulas are
\[
 \begin{gathered}
 \Delta_0(E_i)=E_i\otimes A_i+1\otimes E_i,\qquad
 \Delta_0(F_i)=F_i\otimes1+B_i^{-1}\otimes F_i,\\
 S_0(E_i)=-E_iA_i^{-1},\qquad S_0(F_i)=-B_iF_i.
 \end{gathered}
\]
The square of $S_0$ scales $E_i$ by $q^2$ and $F_i$ by $q^{-2}$ and fixes Cartans. These scalings are automorphisms, so $S_0$ is invertible. Replacing $\Delta_0$ by its opposite coproduct and $S_0$ by $S_0^{-1}$ gives \eqref{eq:hopf}. This operation is called taking the coopposite Hopf algebra; multiplication is unchanged.

By \cite[Proposition 2.73]{VY}, the source Borels possess a pairing $\rho:B_-\times B_+\to\kk$ with
\[
 \rho(N_\lambda,P_\mu)=q^{(\lambda,\mu)},\qquad
 \rho(F_j,E_i)=\delta_{ij}/d
\]
and both mixed root--torus values zero. Its two multiplication rules, using $\Delta_0$, are
\[
 \begin{aligned}
 \rho(yy',x)&=\sum\rho(y,x_{1,0})\rho(y',x_{2,0}),\\
 \rho(y,xx')&=\sum\rho(y_{1,0},x')\rho(y_{2,0},x).
 \end{aligned}
\]
Define $\tau(x,y)=\rho(S_0(y),x)$, as in \cite[Definition 2.74]{VY}. Since $S_0$ reverses multiplication and coproduct, these rules give
\[
 \begin{aligned}
 \tau(xx',y)&=\sum\tau(x,y_{1,0})\tau(x',y_{2,0}),\\
 \tau(x,yy')&=\sum\tau(x_{1,0},y')\tau(x_{2,0},y).
 \end{aligned}
\]
Reversing the coproduct indices gives \eqref{eq:skew}.

For Cartans this construction gives $\tau(P_\lambda,N_\mu)=q^{-(\lambda,\mu)}$. The mixed zero values follow from the corresponding values of $\rho$. Finally,
\[
 \begin{aligned}
 \tau(E_i,F_j)&=-\rho(B_jF_j,E_i)\\
 &=-\rho(B_j,E_i)\rho(F_j,A_i)-\rho(B_j,1)\rho(F_j,E_i)=-\delta_{ij}/d.
 \end{aligned}
\]
Thus existence precedes the calculation of the generator table; it is not inferred from that table alone.
\end{proof}

\subsection{An associative product on the ordered tensor space}
The chosen order of the double is negative followed by positive. We include the construction to specify its multiplication and to justify the tensor-space comparison used later.

\begin{lemma}\label{lem:cocycle}
The vector space $B_-\otimes B_+$ is an associative unital algebra $\cD$ under
\begin{equation}\label{eq:double-product}
 (y\otimes x)(z\otimes u)
 =\sum\tau(x_1,z_1)\tau(x_3,S(z_3))\,yz_2\otimes x_2u.
\end{equation}
Both factors embed as algebras and $(y\otimes1)(1\otimes x)=y\otimes x$.
\end{lemma}
\begin{proof}
Start with the ordinary tensor-product Hopf algebra $C=B_-\otimes B_+$. Define scalar-valued maps
\[
 \begin{aligned}
 \sigma(y\otimes x,z\otimes u)&=\epsilon(y)\tau(x,z)\epsilon(u),\\
 \bar\sigma(y\otimes x,z\otimes u)&=\epsilon(y)\tau(x,S(z))\epsilon(u).
 \end{aligned}
\]
For maps on a coalgebra, convolution is $(f*g)(a)=\sum f(a_1)g(a_2)$. The second rule in \eqref{eq:skew} and the two antipode identities show that $\sigma*\bar\sigma=\bar\sigma*\sigma=\epsilon\otimes\epsilon$. The unit pairing gives $\sigma(a,1)=\sigma(1,a)=\epsilon(a)$.

The required cocycle identity is
\begin{equation}\label{eq:cocycle}
 \sum\sigma(a_1,b_1)\sigma(a_2b_2,c)
 =\sum\sigma(b_1,c_1)\sigma(a,b_2c_2).
\end{equation}
For $a=y\otimes x$, $b=z\otimes u$, $c=w\otimes v$, both sides reduce by \eqref{eq:skew} to
\[
 \epsilon(y)\epsilon(v)\sum\tau(x_1,z)\tau(u,w_1)\tau(x_2,w_2).
\]
Pure tensors span $C$, so the identity holds generally.

Set $a\star b=\sum\sigma(a_1,b_1)a_2b_2\bar\sigma(a_3,b_3)$. To verify associativity explicitly, denote the left and right sides of \eqref{eq:cocycle} by $L(a,b,c)$ and $R(a,b,c)$. Their convolution inverses on $C^{\otimes3}$ are
\[
 \begin{aligned}
 L^{-1}(a,b,c)&=\sum\bar\sigma(a_1b_1,c_1)\bar\sigma(a_2,b_2),\\
 R^{-1}(a,b,c)&=\sum\bar\sigma(a_1,b_1c_1)\bar\sigma(b_2,c_2).
 \end{aligned}
\]
Indeed, each is obtained by reversing the two convolution factors defining $L$ or $R$; ordinary multiplication in $C$ is a coalgebra map. If $m_3(a,b,c)=abc$, substitution and cancellation of adjacent convolution inverses give
\[
 \begin{aligned}
 (a\star b)\star c&=(L*m_3*L^{-1})(a,b,c),\\
 a\star(b\star c)&=(R*m_3*R^{-1})(a,b,c).
 \end{aligned}
\]
For instance the first expression expands as
\[
 \sum\sigma(a_1,b_1)\sigma(a_2b_2,c_1)
 a_3b_3c_2\bar\sigma(a_4b_4,c_3)\bar\sigma(a_5,b_5).
\]
Since $L=R$, associativity follows; normalization gives the unit. Substitution of pure tensors gives \eqref{eq:double-product}. It also gives the asserted factor products. The counits provide linear left inverses to the two factor inclusions.
\end{proof}

\begin{theorem}\label{thm:double}
The identity assignments on $P_r,N_r,E_i,F_i$ give an algebra isomorphism $\cA\cong\cD$.
\end{theorem}
\begin{proof}
The cross product of $x\in B_+$ with $y\in B_-$ is
\begin{equation}\label{eq:cross}
 xy=\sum\tau(x_1,y_1)\tau(x_3,S(y_3))y_2x_2.
\end{equation}
If $c_{ru}=1/n-\delta_{ru}$, then its row and column sums vanish and
\[
 \begin{gathered}
 c_{ri}-c_{r,i+1}=-w_{ri},\\
 c_{ij}-c_{i,j+1}-c_{i+1,j}+c_{i+1,j+1}=-a_{ij}.
 \end{gathered}
\]
Consequently $\tau(P_r,B_i)=\tau(A_i,N_r)=q^{-w_{ri}}$ and $\tau(A_i,B_j)=q^{-a_{ij}}$. Mixed zero values also hold for inverse Cartans, by applying \eqref{eq:skew} to a Cartan times its inverse. In \eqref{eq:cross}, a pair of group-like Cartans has cancelling scalar factors. A Cartan and a root have only one nonzero term. These observations give the first three relations of \eqref{eq:mixed}.

For the remaining relation, the triples in the iterated coproducts are
\[
 (E_i,1,1),\ (A_i,E_i,1),\ (A_i,A_i,E_i)
\]
and
\[
 (F_j,B_j^{-1},B_j^{-1}),\ (1,F_j,B_j^{-1}),\ (1,1,F_j).
\]
Writing $h_{ij}=-\delta_{ij}/d$, their nine contributions to \eqref{eq:cross} are
\[
 \begin{pmatrix}
 h_{ij}B_j^{-1}&0&0\\
 0&F_jE_i&0\\
 0&0&-h_{ij}A_i
 \end{pmatrix}.
\]
The last entry uses $S(F_j)=-F_jB_j$ and $\tau(E_i,F_jB_j)=h_{ij}$. This proves the commutator and defines a homomorphism $\pi:\cA\to\cD$.

Multiplication defines a linear map $\mu:B_-\otimes B_+\to\cA$ using the canonical maps from the factor presentations. It is surjective: count root letters in a word, then count pairs in which a positive letter precedes a negative letter. Swapping an adjacent pair using \eqref{eq:mixed} decreases the second count in the reordered term. The correction to an $E_iF_j$ product replaces two root letters by Cartans, so decreases the first count. Lexicographic induction on these nonnegative counts puts every word in negative--positive order; inverse Cartans obey the same scalar exchange rules.

By Lemma~\ref{lem:cocycle}, $\pi\mu$ is the canonical vector-space bijection onto $\cD$. It follows that $\mu$ is injective. If $\pi(a)=0$, choose $v$ with $\mu(v)=a$; then $\pi\mu(v)=0$ forces $v=0$ and $a=0$. Thus $\pi$ is an isomorphism. No assertion of confluence is needed.
\end{proof}

\section{The full diagonal quotient}\label{sec:quotient}

\begin{proposition}\label{prop:quotient}
In $\cA$, the elements $Z_r=P_rN_r^{-1}$ are central and $\prod_rZ_r=1$. If $J=(Z_r-1:1\leq r\leq n)$, then
\[
 J=(P_r-N_r:1\leq r\leq n),\qquad \cA/J\cong\cT.
\]
The algebra $\cT$ is isomorphic to the full-weight algebra of \cite[Definition 2.13]{VY} at \eqref{eq:base}.
\end{proposition}
\begin{proof}
Cartans commute, and the conjugation weights of $P_r$ and $N_r^{-1}$ cancel on both $E_i$ and $F_i$. This proves centrality. The product-one relations give $\prod Z_r=1$. The identities
\[
 (Z_r-1)N_r=P_r-N_r,\qquad (P_r-N_r)N_r^{-1}=Z_r-1
\]
give equality of ideals. Send both Cartan families in $\cA/J$ to $K_r$, fixing the roots. Every relation becomes a defining relation of $\cT$, including \eqref{eq:single-commutator}. Conversely, send $K_r$ to the class of $P_r$ and fix the roots. The maps are inverse on generators.

For the last assertion send $K_r$ to the source $K_{\lambda_r}$ and fix $E_i,F_i$. The inverse sends $K_\lambda$ to $\prod_rK_r^{m_r}$ for $\lambda=\sum m_r\lambda_r$. Lemma~\ref{lem:lattice} proves independence of the representative. The source $K_{\alpha_i}$ becomes $L_i$, so its commutator, weights and Serre relations coincide with those of $\cT$. This comparison concerns that source-defined algebra; an additional convention identification is needed for the target in \eqref{eq:question-quotient}.
\end{proof}

\begin{proposition}\label{prop:strict}
The ideal $J_{\mathrm{root}}=(A_i-B_i:1\leq i<n)$ is strictly smaller than $J$.
\end{proposition}
\begin{proof}
We have $A_iB_i^{-1}=Z_iZ_{i+1}^{-1}$, so $J_{\mathrm{root}}\subseteq J$. Modulo $J_{\mathrm{root}}$, all $Z_r$ coincide with a central element $z$ satisfying $z^n=1$; this does not force $z=1$. More precisely, there is a homomorphism
\[
 \begin{gathered}
 \cA\longrightarrow\kk[\zeta]/(\zeta^n-1),\\
 P_r\longmapsto\zeta,\quad N_r\longmapsto1,\quad E_i,F_i\longmapsto0.
 \end{gathered}
\]
The element $\zeta$ is a unit, both determinant relations hold, and both $A_i$ and $B_i$ map to one, so every defining relation is respected. The map kills $J_{\mathrm{root}}$ but sends $Z_1-1$ to the nonzero class of $\zeta-1$, whose degree is less than $n$.
\end{proof}

For example, at $n=2$ write $P_1=p=P_2^{-1}$ and $N_1=m=N_2^{-1}$. Then
\[
 [E_1,F_1]=\frac{p^2-m^{-2}}d,
\]
and the full quotient has $p=m=k$ with commutator $(k^2-k^{-2})/d$. The weight Cartan $k$ itself survives; its square is the simple-root Cartan. Although one of the two relations $Z_r=1$ follows from the other, both belong to the specified list of diagonal relations.

\section{Triangular matrices and explicit inverse maps}\label{sec:matrices}

A matrix $T=(t_{ab})$ is a \emph{$q$-quantum matrix} if every submatrix in increasing row and column order,
$\left(\begin{smallmatrix}a&b\\c&e\end{smallmatrix}\right)$, satisfies
\begin{equation}\label{eq:matrix}
 ab=qba,\quad ac=qca,\quad be=qeb,\quad ce=qec,\quad bc=cb,\quad ae-ea=d\,bc.
\end{equation}

\review{This is different than the standard notation of $\left(\begin{smallmatrix}a&b\\c&d\end{smallmatrix}\right)$ and is likely to lead to confusion.}
Its quantum determinant is
\[
 \det_qT=\sum_{\sigma\in S_n}(-q)^{\ell(\sigma)}t_{1,\sigma(1)}\cdots t_{n,\sigma(n)},
\]
where $\ell(\sigma)$ is the number of inversions of the permutation. Let $Q_L$ be the algebra presented by a lower triangular $q$-quantum matrix of determinant one, and $Q_U$ the corresponding upper triangular algebra. Upper or lower forbidden entries are set equal to zero, not removed from the relations before substitution.

In either triangular algebra the diagonal entries commute: a diagonal commutator in \eqref{eq:matrix} contains a forbidden entry. Only the identity permutation survives in the determinant. Indeed, in the lower case a surviving term requires $\sigma(r)\leq r$ for all $r$, and equality of the sums forces equality at every $r$; the upper case is the same with reversed inequalities. Thus if $p_r=t_{rr}$ then
\begin{equation}\label{eq:diag-inverse}
 \prod_rp_r=1,\qquad p_r^{-1}=\prod_{u\ne r}p_u.
\end{equation}
No localization is necessary to invert these entries.

For any root family $X_i$, define interval root vectors recursively by
\begin{equation}\label{eq:root-vectors}
 \begin{aligned}
 R_{b+1,b}(X)&=X_b,\\
 R_{ab}(X)&=R_{a-1,b}(X)X_{a-1}-q^{-1}X_{a-1}R_{a-1,b}(X)\quad(a>b+1).
 \end{aligned}
\end{equation}
They are homogeneous of degree $a-b$ in the simple generators.

\subsection{From entries to simple generators}
\begin{lemma}\label{lem:matrix-to-borel}
There is a surjective homomorphism $\phi_+:B_+\to Q_L$ given by
\begin{equation}\label{eq:phi-positive}
 P_r\longmapsto p_r,\qquad
 E_i\longmapsto d^{-1}p_{i+1}^{-1}t_{i+1,i}.
\end{equation}
For all $a>b$ it satisfies
\begin{equation}\label{eq:root-image}
 \phi_+(R_{ab}(E))=d^{-1}p_a^{-1}t_{ab}.
\end{equation}
\end{lemma}
\begin{proof}
Set $u_i=t_{i+1,i}p_{i+1}^{-1}$. Same-row and same-column relations give
$p_rt_{i+1,i}=q^{w_{ri}}t_{i+1,i}p_r$; other diagonal commutator corrections vanish by triangularity. Hence the proposed image of $E_i$ is $qu_i/d$ and has the required weights.

For adjacent indices, put
\[
 \begin{gathered}
 z=t_{i+1,i},\quad w=t_{i+2,i+1},\quad c=t_{i+2,i},\\
 h=p_{i+1},\quad k=p_{i+2},\quad u=zh^{-1},\quad v=wk^{-1},\quad r=ck^{-1}.
 \end{gathered}
\]
The relevant relations are
\[
 \begin{gathered}
 zw-wz=dhc,\quad zc=qcz,\quad cw=qwc,\quad hc=ch,\quad zh=qhz,\\
 hw=qwh,\quad kz=zk,\quad kw=q^{-1}wk,\quad kc=q^{-1}ck.
 \end{gathered}
\]
Moving the diagonal inverses to the right gives
\[
 quv-vu=dr,\qquad ur=qru,\qquad rv=qvr.
\]
If $R=quv-vu$, the identities $uR=qRu$ and $Rv=qvR$ expand to the two ordered adjacent relations \eqref{eq:serre}, after division by $q$. Distant $u_i,u_j$ commute: their raw commutator correction contains a forbidden upper entry and the relevant diagonals commute with the other entry. Uniform scalar multiplication preserves these homogeneous relations. This proves that \eqref{eq:phi-positive} is a homomorphism.

For $a>b+1$, the matrix relations also give
\begin{equation}\label{eq:entry-recursion}
 t_{a-1,b}t_{a,a-1}-t_{a,a-1}t_{a-1,b}=d\,p_{a-1}t_{ab}.
\end{equation}
Induction on $a-b$, using \eqref{eq:diag-inverse}, proves generation by diagonals and adjacent entries, and hence surjectivity.

For the more precise formula, write $h=p_{a-1}$, $k=p_a$, $x=t_{a-1,b}$, $y=t_{a,a-1}$, $z=t_{ab}$. Besides \eqref{eq:entry-recursion}, we have $kx=xk$ and $yh^{-1}=qh^{-1}y$. Induction in \eqref{eq:root-vectors} gives
\[
 \phi_+(R_{ab}(E))=d^{-2}h^{-1}k^{-1}(xy-yx)=d^{-1}k^{-1}z.
\]
The initial case is \eqref{eq:phi-positive}.
\end{proof}

\subsection{The inverse and its source input}
The root identities needed to reverse this map can be stated without a basis assertion. Write $A=R_{ab}(E)$ and $B=R_{ce}(E)$, with $b<a$, $e<c$ and $(b,a)<_{\mathrm{lex}}(e,c)$. They are
\begin{equation}\label{eq:interval}
 BA=\begin{cases}
 q^{-1}AB,& b=e<a<c\ \text{or}\ b<e<a=c,\\
 AB,&b<e<c<a\ \text{or}\ b<a<e<c,\\
 AB-dR_{cb}(E)R_{ae}(E),&b<e<a<c,\\
 qAB-qR_{cb}(E),&b<a=e<c.
 \end{cases}
\end{equation}
These are precisely the six configurations preceding \cite[Theorem 3.1]{CSS}. To check the normalization, set their parameter to $v=t^{n^2}$ and their rank to $n-1$. Then $v^2=q$ and $v^8\ne1$, as required in their type-A discussion. Their generator $x_{ba}$ is
\[
 x_{ba}=v^{a-b-1}R_{ab}(E),
\]
by their recursive definition. In the crossing case the two product scalings coincide. In the touching case the source correction has one factor $v$ and its long root has one extra factor $v$ relative to the product; together these give the coefficient $v^2=q$. This proves the precise specialization \eqref{eq:interval} of the cited identities in the Serre-presented algebra. It uses neither independence of root monomials nor a triangular-coordinate basis.

\begin{proposition}\label{prop:factor-inverse}
The map $\phi_+$ is an isomorphism, with inverse
\begin{equation}\label{eq:theta-positive}
 t_{aa}\longmapsto P_a,\qquad t_{ab}\longmapsto dP_aR_{ab}(E)\quad(a>b).
\end{equation}
There is also an isomorphism $\phi_-:B_-\to Q_U$. Writing $y_{ab}$ for the entries of $Q_U$ and $h_r=y_{rr}$, it is given by
\begin{equation}\label{eq:phi-negative}
 N_r\longmapsto h_r^{-1},\qquad F_i\longmapsto-d^{-1}y_{i,i+1}h_{i+1}^{-1}.
\end{equation}
Its inverse sends, for $a>b$,
\begin{equation}\label{eq:theta-negative}
 y_{aa}\longmapsto N_a^{-1},\qquad
 y_{ba}\longmapsto(-1)^{a-b}d q^{a-b-1}R_{ab}(F)N_a^{-1}.
\end{equation}
\end{proposition}
\begin{proof}
Let $T_{aa}=P_a$, $T_{ab}=dP_aR_{ab}(E)$ for $a>b$, and $T_{ab}=0$ for $a<b$. Telescoping the simple weights gives
\[
 P_rR_{ab}(E)P_r^{-1}=q^{\delta_{rb}-\delta_{ra}}R_{ab}(E).
\]
Thus the diagonal relations in \eqref{eq:matrix} hold: common columns contribute $q$, common rows contribute $q^{-1}$ when the diagonal is placed first, and disjoint diagonal commutators have a forbidden correction entry.

For completeness consider two distinct strict lower entries with intervals in the order used in \eqref{eq:interval}. A shared start gives
$T_{ab}T_{cb}=qT_{cb}T_{ab}$, since both products have common Cartan factor $d^2P_aP_c$. A shared end gives
$T_{ab}T_{ae}=qT_{ae}T_{ab}$, since both products have common factor $d^2qP_a^2$. Nested intervals give northeast--southwest commutation. Disjoint intervals commute and their northwest--southeast correction is zero, because $T_{ae}=0$.

For crossing intervals $b<e<a<c$, the forward and reverse products have common factor $d^2P_aP_c$. Their difference is
\[
 d^3P_aP_cR_{cb}(E)R_{ae}(E)=dT_{ae}T_{cb}.
\]
Here the two correction roots are nested and commute, and their indicated Cartan reorderings have exponent zero. For touching intervals $b<a=e<c$, the two products are $d^2P_aP_cAB$ and $d^2q^{-1}P_aP_cBA$. Their difference is
\[
 d^2P_aP_cR_{cb}(E)=dT_{aa}T_{cb}.
\]
These are the remaining commutator relations. The six interval configurations exhaust distinct strict lower entries. If a northwest or southeast entry of a rectangle is forbidden, its northeast entry is also forbidden, so its commutator relation is zero on both sides. Relations with forbidden entries in a common row or column are also zero. Finally, the triangular determinant is $\prod P_a=1$. Thus \eqref{eq:theta-positive} defines a homomorphism $\theta_+:Q_L\to B_+$.

Equation~\eqref{eq:root-image} shows that $\phi_+\theta_+$ fixes every matrix entry; the other composite fixes $P_r,E_i$ immediately. This proves the first isomorphism.

Transposition $t_{ab}\mapsto y_{ba}$ preserves the six matrix relations and the triangular determinant. The presentation map
\[
 \chi:B_-\longrightarrow B_+,\qquad N_r\longmapsto P_r^{-1},\quad F_i\longmapsto-q^{-1}E_i
\]
is an isomorphism: its inverse sends $P_r$ to $N_r^{-1}$ and $E_i$ to $-qF_i$, and the weights and homogeneous root relations are respected. Composing $\chi$, $\phi_+$ and transposition gives \eqref{eq:phi-negative}, using
$h_{i+1}^{-1}y_{i,i+1}=qy_{i,i+1}h_{i+1}^{-1}$.

The homogeneity of $R_{ab}$, of degree $m=a-b$, now gives
\[
 \phi_-(R_{ab}(F))=(-1)^m q^{1-m}d^{-1}y_{ba}h_a^{-1},
\]
which proves \eqref{eq:theta-negative}. In particular its adjacent-entry formula is $y_{i,i+1}\mapsto-dF_iN_{i+1}^{-1}$; both the sign and the placement of the Cartan factor are fixed.
\end{proof}

\section{A conditional comparison theorem}\label{sec:comparison}

We now isolate precisely the matrix information needed from a skein algebra. The hypotheses concern an independently given algebra, and include linear independence through a factorization, not just a set of generator relations.

\begin{theorem}\label{thm:comparison}
Let $\cH$ be a unital $\kk$-algebra with a lower triangular $q$-quantum matrix $X=(x_{ij})$ and an upper triangular $q$-quantum matrix $Y=(y_{kl})$, each of determinant one. Assume:
\begin{enumerate}
\item The entry assignments from $Q_L,Q_U$ induce a vector-space bijection
\begin{equation}\label{eq:factorization}
 Q_L\otimes Q_U\longrightarrow\cH,\qquad a\otimes b\longmapsto a(X)b(Y).
\end{equation}
\item For all $1\leq i,j,k,l\leq n$ the mixed exchange identity is
\begin{equation}\label{eq:MX}
\begin{aligned}
 y_{kl}x_{ij}={}&q^{\delta_{ik}-\delta_{jl}}x_{ij}y_{kl}
 +d\,\sgnind{k<i}q^{-\delta_{jl}}x_{kj}y_{il}\\
 &-d\,\sgnind{j<l}q^{\delta_{ik}}x_{il}y_{kj}
 -d^2\sgnind{k<i}\sgnind{j<l}x_{kl}y_{ij}.
\end{aligned}
\end{equation}
Here an indicator $\sgnind{C}$ equals one if $C$ holds and zero otherwise, and forbidden triangular entries are zero.
\end{enumerate}
Then there is an algebra isomorphism $\Phi:\cH\to\cA\cong\cD$ with
\begin{equation}\label{eq:all-images}
 \begin{aligned}
 \Phi(x_{rr})&=P_r,\qquad \Phi(y_{rr})=N_r^{-1},\\
 \Phi(x_{ab})&=dP_aR_{ab}(E)\quad(a>b),\\
 \Phi(y_{ba})&=(-1)^{a-b}dq^{a-b-1}R_{ab}(F)N_a^{-1}\quad(a>b).
 \end{aligned}
\end{equation}
If $I_{\cH}=(x_{rr}y_{rr}-1:1\leq r\leq n)$, then $\Phi(I_{\cH})=J$ and $\cH/I_{\cH}\cong\cT$.
\end{theorem}
\begin{proof}
Put $p_r=x_{rr}$, $h_r=y_{rr}$, $e_i=x_{i+1,i}$ and $f_i=y_{i,i+1}$. By \eqref{eq:diag-inverse}, all $p_r,h_r$ are units. The internal matrix weights are
\[
 p_re_i=q^{w_{ri}}e_ip_r,\qquad h_rf_i=q^{w_{ri}}f_ih_r.
\]
Substitute a diagonal entry into \eqref{eq:MX}. Every correction term contains a forbidden entry of $X$ or $Y$, giving
\begin{equation}\label{eq:diagonal-mixed}
 p_rh_u=h_up_r,\qquad p_rf_i=q^{-w_{ri}}f_ip_r,\qquad
 h_re_j=q^{-w_{rj}}e_jh_r.
\end{equation}
For example, with $x_{ij}=x_{rr}$, the three correction branches contain respectively $x_{kr}$ with $k<r$, $x_{rl}$ with $r<l$, and $x_{kl}$ with $k<r<l$. The case of diagonal $Y$ follows in the same way from its upper triangularity.

For adjacent entries \eqref{eq:MX} becomes
\begin{equation}\label{eq:raw-adjacent}
 f_ie_j=q^{\delta_{i,j+1}-\delta_{i,j-1}}e_jf_i
 +\delta_{ij}d(p_ih_{i+1}-p_{i+1}h_i).
\end{equation}
Indeed the two linear correction branches survive only for $i=j$; the quadratic branch requires $i=j$ and then contains $x_{i,i+1}=0$.

Set $u_j=e_jp_{j+1}^{-1}$ and $v_i=f_ih_{i+1}^{-1}$. The internal weights and \eqref{eq:diagonal-mixed} yield
\[
 \begin{aligned}
 u_jv_i&=q^{\delta_{i,j+1}-\delta_{ij}}e_jf_i p_{j+1}^{-1}h_{i+1}^{-1},\\
 v_iu_j&=q^{\delta_{i,j-1}-\delta_{ij}}f_ie_j p_{j+1}^{-1}h_{i+1}^{-1}.
 \end{aligned}
\]
Thus \eqref{eq:raw-adjacent} implies $[u_j,v_i]=0$ for $i\ne j$, whereas
\[
 [u_i,v_i]=-q^{-1}d(p_ip_{i+1}^{-1}-h_ih_{i+1}^{-1}).
\]
It follows that
\begin{equation}\label{eq:Psi}
 \begin{aligned}
 \Psi(P_r)&=p_r,\qquad\Psi(N_r)=h_r^{-1},\\
 \Psi(E_i)&=d^{-1}p_{i+1}^{-1}e_i=qu_i/d,\\
 \Psi(F_i)&=-d^{-1}f_ih_{i+1}^{-1}=-v_i/d
 \end{aligned}
\end{equation}
satisfy the mixed commutator \eqref{eq:mixed}. Equations~\eqref{eq:diagonal-mixed} give the other mixed relations. Lemma~\ref{lem:matrix-to-borel} applied to $X$ and to $Y^{\mathsf T}$ gives the internal weights and Serre relations. Hence $\Psi:\cA\to\cH$ is a homomorphism.

To prove bijectivity, multiplication $\nu:B_+\otimes B_-\to\cA$ is surjective by the word induction in Theorem~\ref{thm:double}, now exchanging adjacent negative--positive pairs. Solving \eqref{eq:mixed} in that direction still decreases the inversion count or removes two root letters. The composite $\Psi\nu$ is the vector-space bijection \eqref{eq:factorization} transported through the factor isomorphisms of Proposition~\ref{prop:factor-inverse}. If $\nu(v)=0$, this composite forces $v=0$. If $\Psi(a)=0$, choose $v$ with $\nu(v)=a$; again $v=0$. The composite also shows $\Psi$ is surjective. Therefore $\Psi$ is bijective, and its inverse is \eqref{eq:all-images} by the explicit factor inverses. This proves the whole-algebra comparison without presuming that the formulas in the opposite direction respect unspecified additional relations.

Each diagonal generator of $I_{\cH}$ maps to $P_rN_r^{-1}-1=Z_r-1$. Thus $\Phi(I_{\cH})\subseteq J$. Conversely, in any finite sum $\sum a_r'(Z_r-1)b_r'$ representing an element of $J$, lift both coefficients through the surjective map $\Phi$. This gives a preimage in $I_{\cH}$ and proves equality of the full two-sided ideals. Proposition~\ref{prop:quotient} gives the quotient statement.
\end{proof}

\section{Application to the original problem and remaining identifications}\label{sec:application}

Theorem~\ref{thm:comparison} reduces the original question to identifying its hypotheses with the actual skein algebra and its designated arcs. The distinction between a factorization before and after boundary reduction matters. The following elementary lemma supplies that bridge when the required ideal normality is known.

\begin{lemma}\label{lem:quotient-factorization}
Let $H_1,H_2$ be unital subalgebras of an algebra $H$ over a field, such that multiplication $H_1\otimes H_2\to H$ is a vector-space bijection. Suppose $L_i$ is a two-sided ideal of $H_i$ and $HL_i=L_iH$ for $i=1,2$. If $K$ is the ideal they generate in $H$, multiplication induces a vector-space bijection
\[
 (H_1/L_1)\otimes(H_2/L_2)\longrightarrow H/K.
\]
\end{lemma}
\begin{proof}
Normality makes $L_1H$ and $HL_2$ two-sided ideals, so $K=L_1H+HL_2$. Since $H$ is spanned by $H_1H_2$, the internal ideal properties imply
\[
 K=L_1H_2+H_1L_2.
\]
Under the multiplication bijection this is $L_1\otimes H_2+H_1\otimes L_2$. Extending bases of the $L_i$ to bases of the $H_i$ identifies that subspace with the kernel of the tensor quotient map. Taking quotients proves the assertion. The factors need not commute; this is not an isomorphism with an algebra tensor product.
\end{proof}

In \cite{LY}, bad corner arcs have states $i<j$ in the local corner convention; their generated ideal defines the reduced algebra. A punctured surface is essentially bordered when every component has boundary. A disjoint system of boundary-to-boundary arcs is saturated when every complementary component contains exactly one removed point and the system is maximal with this property. For such systems, \cite[Theorem 6.2]{LY} gives a multiplication factorization by arc-neighborhood algebras. Theorem 7.2, equations (128)--(129), gives normality of the corner ideals and expresses the total bad-arc ideal as their sum. These are possible inputs to Lemma~\ref{lem:quotient-factorization}; their application requires distinguishing reduction at one selected corner of a neighborhood bigon from reduction at both corners, and identifying the corresponding triangular quotients.

For the original problem, a sufficient missing geometric statement is the following: the designated $c_{ij}$ are the entries $x_{ij}$ of a lower triangular matrix, the designated $b_{ij}$ are the entries $y_{ij}$ of an upper triangular matrix, and these entries satisfy \eqref{eq:factorization} and the exact identity \eqref{eq:MX}, with the parameter convention \eqref{eq:base}. Rightward orientation alone does not determine these claims. Endpoint state order, framing, stacking, the chosen boundary reduction and the two crossing signs must all be checked. We do not supply that identification here.

If it were established, Theorem~\ref{thm:comparison} would give the prospective original formulas
\[
 \begin{aligned}
 c_{rr}&\mapsto P_r,& b_{rr}&\mapsto N_r^{-1},\\
 c_{ab}&\mapsto dP_aR_{ab}(E)&& (a>b),\\
 b_{ba}&\mapsto(-1)^{a-b}dq^{a-b-1}R_{ab}(F)N_a^{-1}&& (a>b),
 \end{aligned}
\]
with the forbidden entries zero. Since $p_r$ and $h_r$ commute by \eqref{eq:diagonal-mixed}, the products $b_{rr}c_{rr}$ and $x_{rr}y_{rr}$ would coincide under this dictionary. The whole-ideal argument would then apply to every original diagonal generator. These are conditional formulas for the designated arcs, not an asserted map from the original skein algebra.

A second identification is needed on the quantum-group side. Proposition~\ref{prop:pairing} uses the full-weight source presentation, passes to coopposite Borels, and constructs the ordered multiplication \eqref{eq:double-product}. Its generator pairing agrees numerically with \eqref{eq:requested-pairing}. Identifying it with the convention intended in \eqref{eq:question} still requires a dictionary for the original extended Cartans, coproduct, antipode, both skew rules and the order of the double. The comparison of $\cT$ with the source presentation does not by itself provide that dictionary.

Finally, the base in \eqref{eq:base} is an explicit generic choice, not a recovered hypothesis of the original problem. The matrix inverse above invokes the source condition $v^8\ne1$, with $v^2=q$. No conclusion for $q^4=1$, a divided-power form, or a restricted quantum group follows from that argument. A specialization to an algebraic value of $t$ cannot be obtained merely by evaluating a $\Q(t)$-isomorphism: a nonzero polynomial vanishing at that value is invertible in $\Q(t)$. A result over a suitable localized coefficient ring, together with its geometric identifications, is needed for such a deduction.

Thus the pairing normalization, auxiliary inverse maps, injectivity mechanism and full diagonal quotient are available under the stated hypotheses. Establishing their correspondence with the original skein algebra, arcs, coefficients and quantum-group conventions remains necessary for both \eqref{eq:question} and \eqref{eq:question-quotient}.

\begin{thebibliography}{9}
\bibitem{CSS}
Yuqun Chen, Hongshan Shao and K. P. Shum,
\emph{Rosso--Yamane Theorem on PBW basis of $U_q(A_N)$},
arXiv:\href{https://arxiv.org/abs/0804.0954v1}{0804.0954v1} (2008).

\bibitem{LY}
Thang T. Q. L\^e and Tao Yu,
\emph{Quantum traces for $SL_n$-skein algebras},
arXiv:\href{https://arxiv.org/abs/2303.08082v2}{2303.08082v2} (2025).

\bibitem{VY}
Christian Voigt and Robert Yuncken,
\emph{Complex semisimple quantum groups and representation theory},
arXiv:\href{https://arxiv.org/abs/1705.05661v4}{1705.05661v4} (2020).
\end{thebibliography}

\end{document}
